\documentclass[a4paper,12pt]{article}
\usepackage{listings}
\usepackage{xcolor}
\usepackage{hyperref}
\usepackage{geometry}
\usepackage{fontspec}
\setmainfont{Times New Roman}
\geometry{margin=2cm}

\title{Документация к проекту Tetris}
\author{by youngtar}
\date{\today}

% Оформление кода
\lstset{
  basicstyle=\ttfamily\small,
  keywordstyle=\color{blue},
  commentstyle=\color{gray},
  stringstyle=\color{orange},
  breaklines=true,
  frame=single,
  numbers=left,
  numberstyle=\tiny,
  language=C
}

\begin{document}

\maketitle

\tableofcontents

\section{Описание проекта}

Проект реализует консольную игру Tetris на языке C с использованием библиотеки ncurses.

\section{Файловая структура}

\begin{itemize}
  \item \texttt{brick\_game/tetris} — логика игры
  \item \texttt{gui/cli/} — отрисовка
  \item \texttt{tests/} — модульные тесты
  \item \texttt{Makefile} — сборка проекта
  \item \texttt{brick\_game.*} — файлы для запуска игры
\end{itemize}

\section{Как собрать и запустить проект}

\begin{verbatim}
make
make run
\end{verbatim}

\section{Тестирование}

\begin{verbatim}
make test
make gcov_report
\end{verbatim}

\end{document}
